\section{The degree-four stop for the controlled pointwise factor}
\label{sec:136-degree-four-factor-obstruction}

\paragraph{Archive and scope.}
This concerns cote136, physical pp.~117,119,121,
footers116,118,120, scans117,119,121.  It makes the omitted
degree-four case of ARCH-136-LOW-FACTOR-PRESERVATION precise.
The underlying period obstruction is the inherited
ARCH-136-SKELETON-OBSTRUCTION with \(d=3\); no new
counterexample to integral common-source existence is asserted.
We include the exact supplied first skeleton \(F\), rather
than silently replacing it by the first skeleton of the
standard theory.  The conclusion is a failure of unrestricted
preservation under a prescribed degree-four repair.

Let \(B\) be the standard admissible theory and
\(S=\Gamma_{\mathbb Z}(T)\).  Apply the checked
AQ-136-FIRST-SKELETON replacement to obtain an admissible
\(C\) and a strict map
\[
                            \rho:C\longrightarrow B,
                 \qquad C_0=S,\quad C_1=F_1.
                                                               \tag{136.DF1}
\]
Here \(F_1\) is the actual specified free even zero-form
unit-disk algebra, vertically exact in every weight, with its
prescribed two coefficient endpoints.  The left Kan envelope
\(L(F)\) has a canonical map to \(C\); in each simplex
it is graded free and is an \(S\)-linear tensor product of
free even zero-form unit disks.

\paragraph{ARCH-136-DEGREE-FOUR-FACTOR-OBSTRUCTION.}
For every prime \(p\) and all sufficiently large chosen
weights \(m\), there is an old source
\(A\rightrightarrows C\) satisfying all the hypotheses
and the pointwise factor property (136.LF1)--(136.LF2),
with trivial residue grading and literal first skeleton \(F\),
and there is a free graph-relative degree-four Tate attachment
\(A\to E\) with strict target extensions.  For every
choice of those target extensions, \(E\) has nonzero
degree-three graph-difference residues in arbitrarily high
weights.  In particular \(E\) cannot satisfy the same
pointwise factor property for those retained maps.

\begin{proof}
Choose \(m\geq4\) above the fixed comparison bounds for
\(C\) in the finite degrees used below.  The checked
cohomology comparison on the simplicial three-sphere supplies
\[
 \eta\in Z^3 C^m(\Delta[3],\partial\Delta[3]),
 \qquad
                 \int_{\Delta[3]}\rho(\eta)=T^{[m]}.
                                                               \tag{136.DF2}
\]
Indeed both relative third cohomology groups are canonically
\(S^m\), the map \(\rho\) is the identity on coefficients,
and the standard divided-power integration is that canonical
map, with the normalization (136.SO4)--(136.SO6).
Thus an actual closed representative with every face zero
has the asserted period.

Put \(Q=\operatorname{DK}(\mathbb Z[3])\).  Give its
two copies \(V,U\) forms three and four, both of weight
\(m\), and form
\[
 A=L(F)\otimes_{\mathbb Z}
           \bigl(\Lambda(V)\otimes\Gamma(U)\bigr),
                     \qquad dV=pU,\quad dU=0.
                                                               \tag{136.DF3}
\]
Both maps are canonical on \(L(F)\); on the new generators
let \(f_0(V)=f_0(U)=0\), \(f_1(V)=\eta\),
and \(f_1(U)=0\).  The normalized generator has zero
faces, so these are actual strict simplicial differential PD
maps.  The source agrees with \(L(F)\) through simplicial
dimension two, hence has exactly the prescribed \(F\).

In every simplex the algebra (136.DF3) is graded free.
Contract the \(L(F)\) factor \(S\)-linearly.  The first
positive part of the remaining algebra is the injective free
map \(V\xrightarrow{p}U\) in forms three and four;
products of positive generators start in form six.  Therefore
\[
                   H^0(A_r)=S,\qquad
                   H^1(A_r)=H^2(A_r)=H^3(A_r)=0
                                                               \tag{136.DF4}
\]
in every weight.  Rationally every scaled disk is a free unit
disk, so the full rational cohomology is \(S_{\mathbb Q}\).
The weight-zero condition is literal, and \(A_0=S\).
At each prime and simplex take the auxiliary factor to be
\(D=A_{r,(\ell)}\) itself, with identity factor map and
trivial residue grading.  The two evaluations agree on the
graph since they already agree through dimension two.
Equations (136.DF4) verify every part of (136.LF1)--(136.LF2),
including integer torsion-freeness of \(D^0\).

Adjoin a free pair on the graph-relative representable
\(\Delta[3]/\operatorname{sk}_1\Delta[3]\), with both
graph values zero:
\[
              |c|=3,\quad |t|=2,\quad
              \mathrm{wt}(c)=\mathrm{wt}(t)=m,
              \qquad dc=U,\qquad dt=V-pc.
                                                               \tag{136.DF5}
\]
This is the same Tate pair as (136.LF5), after replacing
its generator \(t\) by \(-t\).  Its square-zero differential
and its free pointwise presentation follow directly from
\(dV=pU\).  Since the quotient removes the entire graph,
the resulting \(E\) still has first skeleton \(F\).

Both target extensions exist.  The first sends \(c,t\)
to zero.  For the second take \(c\mapsto0\),
\(t\mapsto\kappa\), where
\[
       \kappa\in C^{m,2}(\Delta[3],\operatorname{sk}_1\Delta[3]),
                                  \qquad d\kappa=\eta.
                                                               \tag{136.DF6}
\]
To see existence, the relative cohomology of
\((\Delta[3],\operatorname{sk}_1\Delta[3])\) with
constant coefficients is concentrated in degree two: the graph
is connected and has reduced cohomology only in degree one.
In the chosen weight range its relative third cohomology in
\(C\) is therefore zero.  This proves (136.DF6) with
actual zero graph restriction, as required by (136.DF5).

Now allow arbitrary extensions \(g_0,g_1:E\rightrightarrows C\)
of the fixed old maps and apply \(\rho\).  Write
\(c_j=\rho g_j(c)\in B_3^{m,3}\),
\(t_j=\rho g_j(t)\in B_3^{m,2}\), and
\(\eta_B=\rho(\eta)\).  Their equations are
\[
                      dt_0=-pc_0,
                      \qquad dt_1=\eta_B-pc_1.
                                                               \tag{136.DF7}
\]
Put \(\partial=d_0-d_1+d_2-d_3\), and in simplicial
degree two define
\[
                   \nu=\partial c\in E_2^{m,3},
                   \qquad b=-\partial t\in E_2^{m,2}.
                                                               \tag{136.DF8}
\]
Every old face of \(U,V\) is zero, so
\(d\nu=0\) and \(db=p\nu\).  Its graph restrictions
are zero because the attachment was relative to the graph.
In simplicial degree two the four new face pairs are independent,
with \(dt_i=-pc_i\), \(dc_i=0\), over \(L(F)_2\).
Contracting that old even-disk factor gives
\(H^2(E_2)=0\) and
\(H^3(E_2)=(S/p)\langle m\rangle^{\oplus4}\).
Thus \([\nu]\) is an actual order-\(p\) class, the
alternating vector of the four face classes.
Moreover \(\rho g_j(\nu)=0\), since \(B_2\) has
no forms of degree three.

Suppose this pair had target primitives with common graph
restrictions in \(C\).  Their images in \(B\) would
give forms \(a_j\in B_2^{m,2}\),
\(h_j\in B_2^{m,1}\), with common graph restrictions and
\[
              da_j=0,\qquad
              dh_j=pa_j-\rho g_j(b)=pa_j+\partial t_j.
                                                               \tag{136.DF9}
\]
Integrate the difference over \(\Delta[2]\) and use
Stokes twice.  Since \(h_1-h_0\) has zero boundary
restriction, one obtains
\[
 \begin{aligned}
 0&=\int_{\partial\Delta[2]}(h_1-h_0)\\
  &=p\int_{\Delta[2]}(a_1-a_0)
                         +\int_{\Delta[2]}\partial(t_1-t_0)\\
  &=p\int_{\Delta[2]}(a_1-a_0)
                         +\int_{\Delta[3]}d(t_1-t_0)\\
  &=T^{[m]}+
       p\left(\int_{\Delta[2]}(a_1-a_0)
                        -\int_{\Delta[3]}(c_1-c_0)\right).
 \end{aligned}                                                \tag{136.DF10}
\]
All displayed periods of integral standard forms lie in \(S^m\).
Equation (136.DF10) is impossible.  More precisely, for any
primitive choices, the boundary period of their difference is
\(T^{[m]}\) modulo \(pS^m\).  This is the nonzero
class in \(H^1(\operatorname{sk}_1\Delta[2];S^m/pS^m)\)
measuring the degree-three graph residue.  Equivalently its
divided linking value is \(T^{[m]}/p\pmod S\).

Multiplying (136.DF8) by \(T^{[w-m]}\), its residue
becomes
\[
                         \binom wm T^{[w]}
                                      \pmod{pS^w}.
                                                               \tag{136.DF11}
\]
For \(w=m+p^a\), \(p^a>m\), this is nonzero and
the weights are unbounded, exactly as in (136.SO7).
It therefore contradicts the uniform high-weight common-graph
primitive consequence of (136.LF1)--(136.LF2), at every
possible proposed cutoff.  Thus \(E\) cannot have that
factor property for any of its retained target extensions.
The weight \(m\) was free to exceed any fixed proposed
degree-four guard from the outset.
\end{proof}

\paragraph{Meaning of the stop.}
The old source already has every low factor required by
LOW-FACTOR-PRESERVATION.  The displayed attachment has actual
target extensions and retains the first skeleton.  Hence that
theorem's excluded degree cannot be restored by assuming only
its stated hypotheses.  This does not say that every carefully
chosen source must contain the old scaled disk in (136.DF3),
or that the complementary degree-four concentration cannot be
completed by retaining further data.  The two targets here
are equal to \(C\), and plainly have a common admissible
source.  The original integral common-source problem remains open.

\paragraph{Verification status.}
The author completed a whole saved-proof reread: exact supplied
first skeleton and incoming replacement; coefficient-normalized
relative representative; full pointwise-factor hypotheses; actual
graph-relative PD attachment and existence of its target maps;
order-p face class; every Stokes sign and arbitrary-choice
term; and unbounded coefficient multiples.  Root and reviewers15 and123 have completed whole saved-proof
readings with PASS.  Reviewer123 contributed to the exact-\(F\)
scope discussion; this is a contributor check, not independent
rediscovery.  The period mechanism is inherited from
the checked SO argument; the exact-\(F\) construction and
its actual degree-four repair are the additional assertions.
